% !TEX root = ../main.tex
%
\chapter{Appendix}
\label{sec:appendix}

\section{Experimental Details}
\label{sec:appendix:exp_details}

Our implementation builds upon the codebase provided by 
\textcite{kuhn2023semanticuncertaintylinguisticinvariances}. 
We utilize the HuggingFace Transformers library to implement 
both the generative model (OPT-6.7B, \citealp{zhang2022opt}) and the 
entailment model (DeBERTa-large-mnli, 
\citealp{he2021debertadecodingenhancedbertdisentangled}).

To ensure a fair comparison, we replicate the generation parameters 
used in the baseline studies. The generation process is divided into 
two distinct phases using the \texttt{generate()} function:

\textbf{Primary Answer Generation:}
To produce the model's ``best guess'' prediction 
(which is evaluated against the ground truth), we employ beam search 
with sampling enabled. Specifically, we set the beam size to 5 
(\texttt{num\_beams=5}) and enable multinomial sampling 
(\texttt{do\_sample=True}).

\textbf{Uncertainty Estimation Sampling:}
To generate the diverse set of sequences required for computing 
semantic entropy and our proposed credal metrics, we disable beam 
search (\texttt{num\_beams=1}) and rely solely on multinomial sampling 
(\texttt{do\_sample=True}). Following the findings of 
\textcite{kuhn2023semanticuncertaintylinguisticinvariances}, 
who demonstrated that uncertainty estimation performance saturates 
beyond a small number of samples, we generate $N=10$ stochastic 
sequences per question.

Detailed implementation details and instructions regarding hardware 
requirements, environment setup,
and the exact commands required to reproduce the tables and figures
presented in Section 5 are provided in the repository (
\url{https://github.com/liamhpr/Credal-Prediction-for-LLMs.git}).

The experiments run on one 80GB NVIDIA A100.



For our evaluation, we split the temperature interval $[0.4, 2.2]$ using
20 and 50 steps. The step size stays consistent and the resulting 
temperature spaces can be seen in the two Tables 
\ref{tab:twenty_steps_interval}, \ref{tab:fifty_steps_interval} below.

Table \ref{tab:twenty_steps_interval}:
The temperature interval $[0.4, 2.2]$ split into 20 evenly spaced steps. The
temperature yielding the maximum average likelihood ($\tau^{ML} = 1.347$) is
highlighted in bold. In addition, the resulting $\alpha$-cuts
$\mathcal{T}_\alpha$ with $\alpha \in \{0.6, 0.7, 0.8, 0.9\}$ are shown.
\begin{table}[h]
    \centering
    \begin{tabular}{p{0.85\textwidth}}
    \toprule
    \textbf{Set of Evaluated Temperatures ($\mathcal{T}$)} \\
    \midrule
    0.400, 0.495, 0.590, 0.684, 0.779, 0.874, 0.968, 1.063, 1.158, 
    1.253, \\
    \textbf{1.347}, 1.442, 1.537, 1.632, 1.726, 1.821, 1.916, 2.011, 2.105, 2.200 \\
    \midrule
    $\mathcal{T}_{0.9} = \{ 1.253, \dots, 1.537\}$ \quad \quad
    $|\mathcal{T}_{0.9}| = 4$ \\
    $\mathcal{T}_{0.8} = \{ 1.158, \dots, 1.632 \}$ \quad \quad
    $|\mathcal{T}_{0.8}| = 6$ \\
    $\mathcal{T}_{0.7} = \{ 1.063, \dots, 1.726 \}$ \quad \quad
    $|\mathcal{T}_{0.7}| = 10$ \\
    $\mathcal{T}_{0.6} = \{ 1.063, \dots, 1.821\}$ \quad \quad
    $|\mathcal{T}_{0.6}| = 11$ \\
    \bottomrule
    \end{tabular}
    \caption{}
    \label{tab:twenty_steps_interval}
\end{table}

Table \ref{tab:fifty_steps_interval}:
The temperature interval $[0.4, 2.2]$ split into 50 evenly spaced steps. The
temperature yielding the maximum average likelihood ($\tau^{ML} = 1.392$) is
highlighted in bold. In addition, the resulting $\alpha$-cuts
$\mathcal{T}_\alpha$ with $\alpha=0.9$ are shown.
\begin{table}[h]
    \centering
    \begin{tabular}{p{0.9\textwidth}}
    \toprule
    \textbf{Set of Evaluated Temperatures ($\mathcal{T}$)} \\
    \midrule
    0.400, 0.437, 0.474, 0.510, 0.547, 0.584, 0.620, 0.657, 0.694, 0.731, 
    0.767, 0.804, 0.841, 0.878, 0.914, 0.951, 0.988, 1.025, 1.061, 1.098, 
    1.135, 1.171, 1.208, 1.245, 1.282, 1.318, 1.355, \textbf{1.392}, 1.429, 1.465, 
    1.502, 1.539, 1.576, 1.612, 1.649, 1.686, 1.723, 1.759, 1.796, 1.833, 
    1.869, 1.906, 1.943, 1.980, 2.016, 2.053, 2.090, 2.127, 2.163, 2.200 \\
    \midrule
    $\mathcal{T}_{0.9} = \{ 1.208, \dots, 1.576 \}$ \quad \quad
    $|\mathcal{T}_{0.9}| = 11$\\
    \bottomrule
    \end{tabular}
    \caption{}
    \label{tab:fifty_steps_interval}
\end{table}
